%% Appendix C -- calibration-block length and duty cycle (Phase 3, 2026-09-28): moved verbatim from Sect. 6.2,
%% with three self-references adjusted. Numbers unchanged.
\section{Calibration-block length and duty cycle}
\label{app:C}

This appendix gives the measurements behind the second paragraph of Sect.~\ref{sec:6.2}. They are not conclusive, for the
reason stated below, and are recorded so that the measurement that would settle the question can be designed.

What acts on knot noise is the number of scans averaged into each
calibration block. Splitting each block into contiguous quarters separates the
two components that limit a block mean -- white scan noise $\sigma_s$, which
averages down, and drift within the block at rate $g$, which does not:
\begin{equation}
\begin{split}
\mathrm{scale}(Q_0-Q_1)^2 &= \frac{8\sigma_s^2}{N} + \Bigl(\frac{gT}{4}\Bigr)^2,\\
\mathrm{scale}\Bigl(\tfrac{Q_0+Q_1}{2}-\tfrac{Q_2+Q_3}{2}\Bigr)^2
 &= \frac{4\sigma_s^2}{N} + \Bigl(\frac{gT}{2}\Bigr)^2 .
\end{split}
\label{eq:18}
\end{equation}

Sub-blocks must be cut contiguously for this, because an interleaved
(even/odd) split distributes neighbouring scans across both halves and
cancels whatever they share. The interleaved split is kept as a control: the
ratio between the two, at the same number of scans, measures how much of the
block's scatter is common to neighbouring scans rather than independent.

The operational block is $N = 34$ scans spanning $T = 32.0$\,s, in both
channels and for both reference gases. Within it, averaging does not behave as
white noise (Table~\ref{tab:subblock_scaling}).

\begin{table}[t]
\caption{Scatter of contiguous sub-block contrasts within the operational calibration block ($N = 34$, $T = 32.0$\,s) in the two heated channels, the averaging exponent they imply, and the ratio of contiguous to interleaved scatter at the same number of scans $m$, with the values expected for white noise.}
\label{tab:subblock_scaling}
\begin{tabular}{llll}
\tophline
 & hot ANs & hot PNs & white noise \\
\middlehline
contiguous quarters (8 scans) & $\fieldnum{2.506\times10^{-4}}$ & $\fieldnum{1.095\times10^{-4}}$ & \\
contiguous halves (16 scans) & $\fieldnum{2.332\times10^{-4}}$ & $\fieldnum{0.820\times10^{-4}}$ & \\
\textbf{averaging exponent} & \textbf{\fieldnum{0.104}} & \textbf{\fieldnum{0.416}} & 0.500 \\
contiguous / interleaved, same $m$ & \textbf{\fieldnum{2.26}} & \textbf{\fieldnum{1.22}} & 1.03 \\
\bottomhline
\end{tabular}
\end{table}

Doubling the number of scans averaged reduces the scatter by \fieldnum{7}\,\% in the ANs
channel, against the 29\,\% that independent scans would give. Most of what a
block mean contains is therefore not independent between scans, and the last
row says so directly: cutting the same block contiguously rather than in
alternation more than doubles its apparent scatter in ANs. These are spectral-domain statistics: in the
NO$_2$ amount, contiguous and interleaved halves give similar knot noise (Sect.~\ref{sec:4.2}), so most of this
correlated structure is evidently absorbed by the baseline polynomial before it reaches the gas columns. A monotonic trend
across the 32\,s is present in both channels but explains little of this -- see
below.

Two conclusions follow, and the second is a limitation. Within the range that
can be measured, lengthening the averaging window is worth much less than the
$N^{-1/2}$ a reader would assume, and how much less depends on the channel.
And the range that can be measured is entirely below the operational block
length: subdividing $N$ scans says nothing about $2N$ or $4N$. The evidence
available is therefore not neutral but mildly unfavourable, and the
measurement that would settle it -- operating one longer block -- has not been
made.

The duty cycle is not what stands in the way. Calibration occupies 1.2\,\% of
the measurement time (Sect.~\ref{sec:6.4}), and a fourfold zero-air block would take it to
3.9\,\% (Table~\ref{tab:duty_blocks}).

\begin{table}[t]
\caption{Calibration duty cycle for the operational zero-air block and for blocks two and four times longer.}
\label{tab:duty_blocks}
\begin{tabular}{llll}
\tophline
zero-air block & $N$ & zero-air duty & total calibration duty \\
\middlehline
operational & 34 & 0.90\,\% & 1.20\,\% \\
$\times 2$ & 68 & 1.79\,\% & 2.10\,\% \\
$\times 4$ & 136 & 3.58\,\% & 3.89\,\% \\
\bottomhline
\end{tabular}
\end{table}

Even under the white-noise assumption the return would be modest, and it is
worth stating why before anyone spends duty cycle on it. Halving the zero-air
knot noise removes three quarters of the $I_0$ variance, but that term carries
\fieldnum{45}\,\% of the $\Sigma$ANs variance, so the structural budget would fall from
\fieldnum{0.1236} to \fieldnum{0.1008}\,ppb -- \fieldnum{18.5}\,\%, not half. The measured exponents put even that
out of reach in one channel; Table~\ref{tab:ceilings} gives the ceiling for each remedy.

\begin{table}[t]
\caption{Best attainable $\Sigma$ANs structural uncertainty (standard deviation) under each remedy, as a fraction of the operational value, and the assumption each row rests on. These are ceilings, not forecasts.}
\label{tab:ceilings}
\begin{tabular}{l>{\raggedright\arraybackslash}p{0.28\textwidth}l>{\raggedright\arraybackslash}p{0.3\textwidth}}
\tophline
 & $\Sigma$ANs structural (SD), best attainable & of operational & status \\
\middlehline
operational & \fieldnum{0.1236}\,ppb & 1.00 & measured \\
$\times 4$ blocks only & \fieldnum{0.1008} & \fieldnum{0.82} & assumes $N^{-1/2}$; not supported below $N$ \\
flag bad $R$ intervals only & \fieldnum{0.0901} & \fieldnum{0.73} & assumes a perfect detector \\
\textbf{both} & \textbf{\fieldnum{0.0547}} & \textbf{\fieldnum{0.44}} & both assumptions \\
\bottomhline
\end{tabular}
\end{table}

These are ceilings, not forecasts, and the way they are computed needs stating
because Sect.~\ref{sec:4.3} forbids the obvious method. Removing a term's contribution from a
total is a quadrature operation, and the joint propagation is not a quadrature
sum -- here it exceeds one by \fieldnum{7.6}\,\%, \fieldnum{0.1236}\,ppb against \fieldnum{0.1149}. The rows are
therefore evaluated in quadrature space and rescaled by that measured ratio,
which makes the operational row reproduce the measured joint exactly and
applies the same correction to the others. The residual assumption is that the
ratio survives shrinking a term, which it will not do exactly, because the
ratio depends on the covariance and interaction of the terms (Sect.~\ref{sec:4.3}), and shrinking one
term changes them.

The ratio column is unaffected by any of this: the rescaling factor cancels,
so the relative ordering is the same whether the calculation is done in
quadrature space or after rescaling. Section~\ref{sec:6.2} uses only the ordering.

Flagging outranks averaging here, and it does so under every basis we have
checked -- record space as well as hourly, and quadrature space as well as
rescaled. The flagging row assumes a detector that
removes the tail and nothing else, reducing the $R$ term to its present robust
scale; a real flag will fall short. The block row assumes an averaging
behaviour the measurements of this appendix do not support.

A linear trend across the block is part of the explanation but not all of it.
It accounts for \fieldnum{40}\,\% of the within-block variance in the ANs channel and \fieldnum{34}\,\%
in PNs, inflating the apparent sub-block scatter by \fieldnum{1.29} and \fieldnum{1.23}. Whether
removing it restores white averaging cannot be tested by detrending and
re-measuring: the contrast comparing contiguous halves and the least-squares
slope over four sub-blocks have a cosine of $-0.894$, so subtracting the line
subtracts most of the quantity being measured. Measured instead on a contrast
orthogonal to the trend, the residual scatter is still \fieldnum{1.88} times white in ANs
and \fieldnum{1.17} in PNs. Detrending is a partial remedy; what limits averaging is a
correlation that outlives the removal of its linear part.

What does not depend on either assumption is the ordering. The reflectivity
term is concentrated in a few episodes and does not average down; the zero-air
term is distributed and averages at least partly. Flagging comes first
whatever the detector achieves and whatever the block scaling turns out to be.

One observation from the same measurement is worth recording. Both channels
show a monotonic trend across the 32\,s block, of opposite sign -- \fieldnum{$+248$} to
\fieldnum{$-314$}\,ppm in ANs against \fieldnum{$-261$} to \fieldnum{$+235$}\,ppm in PNs. Two adjacent windows
moving in opposite directions is a change of spectral slope rather than of
intensity, which points at the source or at the bandpass separating the
channels rather than at the gas fill. Removing the linear trend by contrast
does not reduce the excess scatter, so this trend is not what limits
averaging; the correlation between neighbouring scans is.
